\section{Implementation}

\subsection{Implementation Environment}

The implementation is based on gem5 version 23.0.0.1 and is developed
on a separate branch.\footnote{Source code: \url{https://github.com/ayushgangal/gem5-UGRC-} (branch \texttt{astra-vastra-dev}).} The simulator is built for the RISC-V ISA as the
\path{build/RISCV/gem5.opt} target. All experiments use
syscall-emulation mode with a single \texttt{RiscvMinorCPU} that has
one hardware thread. The workloads are freestanding, statically linked
RISC-V programs built for the \texttt{rv64g} ISA, so compressed
instructions are not used.

The baseline configuration, \path{configs/research/vp_baseline.py},
models a 1\,GHz \texttt{RiscvMinorCPU} with timing memory accesses,
64-byte cache lines, a 32\,KiB 4-way L1 instruction cache, a 64\,KiB
4-way L1 data cache with 16 MSHRs and a stride prefetcher, a 256\,KiB
8-way L2 cache with 20 MSHRs and an AMPM prefetcher, and DDR4-2400
main memory. The pipeline widths are four, at most one memory
instruction is issued and one committed per cycle, early memory issue
is enabled, and the CPU uses the \texttt{SingleThreaded} thread policy
with up to 16 outstanding memory accesses.

\subsection{LVP Integration}
\label{sec:lvp-integration}

The LVP predictor is integrated into the Execute stage of gem5
MinorCPU. The integration involves more than looking up a predicted
value: Execute must associate each prediction with the correct dynamic
producer, determine when a dependent load may use it, substitute the
predicted operand while generating the memory request, and validate
the prediction when the producer's actual response becomes available.
Incorrect predictions that have influenced speculative execution
require the younger instruction stream to be discarded and fetched
again.

The implementation described here covers standalone LVP-based
load-to-load address speculation. The arbitration between LVP and
V-ASTRA when META is enabled is described separately in
Section~\ref{sec:impl-meta}.


\subsubsection{Configuration and Prediction Eligibility}

Three configuration conditions determine whether LVP-based address speculation can operate. \texttt{enableLVP} enables the predictor, while
\texttt{lvpSpecAddr} enables dependent loads to use confident
predictions as their base addresses. Enabling address speculation
requires \texttt{enableLVP}, \texttt{lvpSpecAddr}, and
\texttt{executeAllowEarlyMemoryIssue} to be enabled. Invalid
combinations are rejected during construction of the Execute stage.

An optional PC range can further restrict address speculation.
When configured, both the producer and consumer instruction PCs must
belong to the half-open interval
\[
    [\,\texttt{lvpSpecPcBegin},\ \texttt{lvpSpecPcEnd}\,).
\]
Setting both bounds to zero leaves the range unrestricted.

By default, the predictor uses a 1,024-entry, untagged table
indexed by the load instruction's PC:

\[
    \operatorname{index}(PC)
    = (PC \gg 1) \mathbin{\&} (N-1),
    \qquad N=1024.
\]

Each entry stores the most recently observed value and a 2-bit
saturating confidence counter. The counter has four states, from
zero to three, and the default prediction threshold is three (it
can be configured from one to three). A lookup returns a prediction
only when the entry's confidence reaches the configured threshold.
On a matching observation, the counter is incremented up to its
maximum. On a mismatch, it is decremented down to zero. The stored
value is replaced by the newly observed value in either case.
Entries initially contain value zero and confidence zero. The table
does not maintain a separate valid bit or PC tag.

The implementation also restricts which instructions can act as
producers and consumers. An eligible producer is an ordinary integer
load with one integer destination register, excluding the zero
register. Atomics, store-conditionals, micro-operations, serializing
instructions, and prefetches are excluded. An eligible consumer must
satisfy the same ordinary-load restrictions and have exactly one
integer source register, corresponding to its base address.

Consequently, the implemented speculative operand path is limited to
dependent loads. General ALU instructions do not directly consume
LVP predictions.


\subsubsection{Issue-Time Prediction and Dependency Tracking}

Prediction lookup and dependency tracking are performed in
\texttt{Execute::lvpOnIssue()}, which is called when an instruction
is issued into a functional unit. The function is invoked in
instruction issue order.

The predictor table and \texttt{lvpSpecMap} serve different purposes.
The predictor table stores learned values indexed by static
instruction PCs. In contrast, \texttt{lvpSpecMap} tracks the
currently available predicted writer for each architectural register.
It is maintained separately for each thread and maps a flattened
destination register to an entry containing the producer's dynamic
sequence number, PC, and predicted value:

\[
\texttt{lvpSpecMap[tid][reg]}
=
\{\texttt{producerSeq},\texttt{producerPC},\texttt{value}\}.
\]

The issue-time processing consists of three ordered operations.
The predictor lookup in the third step is performed whenever LVP is
enabled. The consumer lookup, destination invalidation, and insertion
into \texttt{lvpSpecMap} are performed only when address speculation
(\texttt{lvpSpecAddr}) is also enabled. Without it, predictions are
made and later validated, but no consumer can use them.

\paragraph{1. Consumer dependency lookup.}
If the issued instruction is an eligible consumer, Execute examines
its base register and searches for a corresponding entry in
\texttt{lvpSpecMap}. Finding an entry is not sufficient. The recorded
producer sequence number must equal the consumer's
\texttt{instToWaitFor} value, which identifies the dynamic instruction
on which the consumer is waiting. Both the consumer PC and the
recorded producer PC must also satisfy the configured speculation
window.

When these checks succeed, the consumer copies the producer identity,
source operand index, source register, and predicted value into its
own dynamic instruction metadata. This snapshot preserves the
specific producer-consumer relationship even if the map is updated
later.

\paragraph{2. Destination invalidation.}
Execute then removes existing map entries for every destination
register written by the current instruction. This prevents a
prediction belonging to an older writer from remaining available
after a newer instruction has become the latest writer of that
register. This invalidation is not restricted to eligible loads:
any instruction that writes a register invalidates the corresponding
entry. The no-cost instruction path also performs destination
invalidation.

The ordering of the first two operations is significant. For an
instruction such as \texttt{ld a0, 0(a0)}, the source register must
first be checked against the previous predicted writer. Only
after the consumer has captured that dependency can the instruction's
own destination write invalidate the older mapping.

\paragraph{3. Producer prediction lookup.}
Finally, if the current instruction is an eligible producer,
Execute queries the LVP table using its PC. If the lookup returns a
confident prediction, the dynamic instruction is marked as predicted,
its predicted value and producer state are recorded, and, when
address speculation is enabled, the prediction is inserted into
\texttt{lvpSpecMap} under its flattened destination register. The
entry identifies this particular dynamic instance using its execution
sequence number.

The sequence-number association is essential because multiple
dynamic instructions can write the same architectural register.
A consumer must use the prediction belonging to the exact producer
identified by its dependency information, rather than an arbitrary
prediction associated with that register.

\subsubsection{Admission for Early Memory Issue}

Saving a prediction in the consumer's metadata does not immediately
cause speculative execution. The consumer must subsequently reach
the end of its memory functional unit and pass the admission checks
in \texttt{Execute::lvpSpecAdmissible()}.

MinorCPU maintains an in-flight instruction queue and a separate
queue of memory references waiting to leave their functional units
for the LSQ. The normal early-memory-issue path considers the
oldest unissued memory reference, checks that it has reached the
front of its functional unit, and verifies that the LSQ can accept
another request.

Under normal execution, a memory instruction whose dependency is
still outstanding cannot issue its request. LVP introduces an
additional admission path for the case in which the unresolved
dependency is the predicted load at the head of the in-flight
queue.

The caller attempts LVP admission only when the consumer's
dependency has not yet been satisfied through normal execution.
If the dependency has already been satisfied, the instruction can
proceed through the ordinary early-memory-issue path without
requiring LVP speculation.

The LVP-specific admission function enforces the following
conditions:

\begin{enumerate}
    \item LVP and address speculation must both be enabled.

    \item The candidate must be the oldest memory reference waiting
    to enter the LSQ, and it must belong to the same thread as the
    in-flight head instruction.

    \item The consumer must have a saved speculative operand and
    producer identity. Its \texttt{inst\allowbreak ToWaitFor} value must still
    match the saved producer sequence number, and its consumer
    state must still be \texttt{None}.

    \item The in-flight head must be the predicted producer identified
    by the consumer. That producer must still have state
    \texttt{Predicted}, must already be in the LSQ, and must not
    have a memory response available.

    \item Both the producer and consumer must belong to the current
    instruction stream.
\end{enumerate}

The head-instruction requirement is particularly important. It
ensures that the predicted producer is the oldest outstanding
instruction. Older instructions have therefore already left the
in-flight queue, and the candidate consumer is not bypassing an
arbitrary chain of unresolved dependencies. Requiring the producer
to be in the LSQ but not yet have a response also ensures that
speculation is attempted only while the actual load result is still
pending.

The admission function is only the LVP-specific part of the
early-issue decision. The surrounding Execute logic continues to
enforce structural conditions, including LSQ availability, functional
unit completion, and the requirement that the candidate has not
already entered the LSQ.


\subsubsection{Operand Substitution and LSQ Request Generation}

After admission succeeds, Execute temporarily enables operand
substitution using \texttt{LVPSpec\allowbreak OperandScope}. This scope is
constructed around the call to \texttt{commitInst()} and activates
the consumer's \texttt{lvpSpec\allowbreak Active} flag. Its destructor clears
the flag when the scope ends, including when control exits through
an early return.

The actual substitution is implemented in
\texttt{ExecContext::\allowbreak getRegOperand()} in
\texttt{exec\_\allowbreak context.hh}.
When an instruction requests a source operand, the function returns
the saved predicted value only when all of the following conditions
hold:

\begin{enumerate}
    \item Speculative operand substitution is currently active.
    \item The requested source index matches the saved source index.
    \item The static instruction requesting the operand is the same
    static instruction associated with the dynamic instruction.
    \item The requested register matches the saved source register.
\end{enumerate}

If any of these checks fails for a valid source register, the
function returns the normal architectural register value.
The substitution therefore applies only to the specific base-register
operand of the admitted consumer. It does not modify the architectural
register file.

During \texttt{commitInst()}, the load's normal
\texttt{initiateAcc()} path obtains the substituted base operand
and forms the effective address. The resulting memory request is
sent through the ordinary memory-access path into the LSQ. Thus,
the consumer issues a real load request using the predicted address.
The request is not merely an address calculation or a simulated
access.

The consumer is marked as
\texttt{LVPConsumerState::Accepted} only after the request has
actually entered the LSQ. Execute records this transition and
updates the corresponding consumer statistic. After
\texttt{commitInst()} returns, the producer's
\texttt{lvpUsedBySpec} flag is set only if the consumer reached
the \texttt{Accepted} state.

This distinction separates three events: saving a candidate
prediction, passing admission checks, and actually issuing a
memory request using the prediction. Only the last event counts
as speculative use and can make a later producer misprediction
require recovery.

\subsubsection{Producer Validation and Misprediction Recovery}

Validation takes place when the producer's memory response is
handled. Since Execute processes completed instructions in
in-flight order, the producer reaches the head before a younger
consumer can complete architecturally.

For a successful load response, \texttt{completeAcc()} processes
the returned data and writes the actual loaded value into the
architectural destination register. Execute then calls
\texttt{lvpValidate\allowbreak Producer()} for an eligible load. The function
reads the actual destination-register value and trains the LVP
table using the producer PC and this observed value. Training is
performed for successful eligible loads regardless of whether
the load previously received a confident prediction.

If the producer was predicted, the actual value is compared with
the saved prediction and the validation outcome is recorded as
either correct or incorrect. A recovery is requested only when
the prediction was incorrect and the producer's
\texttt{lvpUsedBySpec} flag is set. An incorrect prediction that
was never used by an accepted speculative consumer therefore
does not trigger a redirect.

When recovery is required, Execute first completes the producer's
response handling and processes its ordinary control-flow outcome.
It then generates an \texttt{UnpredictedBranch} redirect to the
thread's current PC, which is the instruction following the
completed producer.

The producer is retained. The redirect discards all younger
instructions belonging to the old stream, not merely those
that directly depended on the incorrect prediction. These
instructions are subsequently fetched and executed again
from the producer's next PC. The redirect is guarded by a check
that there is not already a stream change in the current cycle.

A stream change increments the thread's stream sequence number
and clears its \texttt{lvpSpecMap}. Younger instructions from the
previous stream are subsequently discarded through MinorCPU's
normal stream-mismatch handling. This also prevents a speculative
memory response from an obsolete stream from being completed as
a valid instruction.

As an additional safeguard, when an accepted speculative consumer's
own memory response is handled, the implementation checks that the
architectural value of its base register equals the predicted value
it used. Because the producer has committed by then, the two must be
equal. A mismatch stops the simulation. This check therefore exposes
any failure of the recovery path instead of allowing an incorrectly
addressed load to commit.

An additional Fetch2 modification handles the branch-predictor
state associated with this redirect. Unlike an ordinary branch
misprediction, an LVP recovery originates from a load instruction,
which was not itself predicted by the branch predictor. However,
younger branch instructions may already have updated speculative
predictor history, return-address-stack state, or indirect-predictor
state. The \texttt{lvpRedirect} flag identifies this special case
in \texttt{Fetch2::\allowbreak updateBranch\allowbreak Prediction()}. Fetch2 clears the
flag, invokes the branch predictor's \texttt{squash()} and
\texttt{update()} operations using the redirecting instruction's
fetch sequence number and thread ID, and returns without applying
the ordinary branch-prediction handling to the load.

The Execute redirect and the Fetch2 predictor update therefore
serve different purposes: the former discards the incorrect
instruction stream, while the latter restores the branch predictor's
speculative bookkeeping consistently with that redirect.

Each dynamic instruction tracks its producer and consumer roles
separately, so a dependent load can consume an older prediction while
its own value is also predicted. When an instruction leaves the
in-flight queue, \texttt{lvpClassify()} records whether each role
ended squashed, faulted, or committed. A producer's
\texttt{lvpSpecMap} entry is removed only if it still identifies that
same dynamic producer, so an older producer's cleanup cannot delete a
newer writer's prediction.

\subsection{ASTRA Address Prediction}
\label{sec:impl-astra}

ASTRA is implemented as a PC-indexed predictor consisting of a
Last Address Table (LAT) and a Stride$\rightarrow$Stride table (SS). The
predictor is implemented in \texttt{astra.hh} and \texttt{astra.cc},
with prediction lookup and training integrated into the Execute stage
of MinorCPU.

Unlike the paper's plain ASTRA, whose LAT holds the last committed
virtual address of a load, this implementation records the load's
last base-register value in the LAT and predicts the next base value,
as in the enhanced V-ASTRA design. The
instruction's displacement is added only when V-ASTRA forms an
effective address (Section~\ref{sec:impl-vcache}). Addresses are kept
at their full 64-bit width.

\subsubsection{Predictor Organization}

Each LAT entry contains two fields: the most recently trained
base-register value and an index into the SS table.
The base is stored as a 64-bit \texttt{Addr}, while the SS contains
signed 16-bit stride values. Both tables contain
\texttt{astraTableEntries} entries, with a default size of 1,024.

The LAT row \(r\) of a load is selected using its PC:

\[
    r = (PC \gg 1) \mathbin{\&} (N-1),
\]

where \(N\) is the configured number of entries. A stride \(S\)
selects SS entry \(S \bmod N\), that is, its lower bits.

The two tables are untagged and do not maintain separate valid
bits or confidence counters. Initially, all LAT entries contain
a base of zero and an SS index of zero, while all SS entries
contain a stride of zero. Consequently, every eligible lookup
produces a prediction, including during the initial cold state.

\subsubsection{Prediction and Base Capture}

Prediction lookup is performed by \texttt{Execute::astraOnIssue()},
which is called from the instruction issue path. The function first
checks whether the instruction is eligible for ASTRA prediction.
Eligible instructions are ordinary loads that have an integer
base-register source. Atomic operations, store-conditionals,
micro-operations, serializing instructions, and data prefetches
are excluded. Both floating-point loads and compressed loads
can be looked up by ASTRA because their base address is supplied
through an integer register.

For an eligible load, the predictor uses its LAT entry to obtain
the most recently stored base, \(L\), and the SS index stored with
it, \(k\). The indexed stride is
sign-extended from 16 bits and added to the stored base, using 64-bit
arithmetic without overflow detection:

\[
    B = L + SS[k].
\]

The predicted base \(B\) is stored in the dynamic instruction's
\texttt{astraPredBase} field, and its ASTRA state is marked as
\texttt{Looked}.

The base that the load actually uses is captured separately by
\texttt{Execute::\allowbreak astraCapture\allowbreak Base()} in
\texttt{executeMemRefInst()}, immediately before the load's
\texttt{initiateAcc()} call, and stored in \texttt{astraActualBase},
with \texttt{astraBaseValid} indicating that a usable base was
captured. The captured value depends on how the load obtains its base
operand:

\begin{itemize}
    \item Ordinarily, it is the architectural value of the base
    register, which in MinorCPU holds the committed value at this
    point.

    \item If the load's base was supplied by a validated V-ASTRA value
    through dependent-load forwarding (Section~\ref{sec:impl-sb}), the
    forwarded value is captured. It is the value the producer will
    write, but the producer has not yet committed.

    \item If the base was supplied by an LVP prediction, it is not
    captured in the default configuration, and
    \texttt{astraBaseValid} remains false. In the experimental modes of
    Section~\ref{sec:impl-provisional}, the LVP-supplied value is
    captured as a provisional, not yet validated, base.
\end{itemize}

In standalone ASTRA operation, the prediction is used for
measurement and comparison rather than directly controlling
memory execution. ASTRA itself does not calculate the effective
address or suppress the ordinary load access. The use of the
predicted base to form a candidate effective address belongs
to the V-ASTRA mechanism described in Section~\ref{sec:impl-vcache}.

\subsubsection{Training and Predictor Lifecycle}

ASTRA training occurs when an instruction reaches the head of
the in-flight queue and retires. The Execute commit path invokes
\texttt{Execute::astraRetire()} immediately before removing
the instruction from the in-flight queue. Therefore, prediction
lookup occurs at issue, while the corresponding predictor update
takes place at retirement.

For a load with a valid captured base \(A\), the predictor
calculates the stride relative to the base \(L\) stored in its
LAT entry:

\[
    S = A - L.
\]

The stride is represented using 16 bits. Consequently, differences
outside the signed 16-bit range wrap when represented by the
predictor.

The SS is updated using the index \(k\) stored in the LAT entry
before the update:

\[
    SS[k] \leftarrow S.
\]

The LAT is then updated with the newly observed base and
the index associated with the new stride:

\[
    LAT[r] \leftarrow
    \bigl(A,\ S \bmod N\bigr).
\]

This update order allows the predictor to learn transitions
between successive strides. The SS entry selected by the
previous stride context is updated with the newly observed
stride, after which the LAT records the context for the next
prediction.

Only successfully committed loads with a valid captured base
train the predictor. Faulting and squashed instructions are
classified separately and do not update the tables. A committed
instruction for which the base is unavailable is also classified
without training.

Since prediction occurs at issue but training occurs at retirement,
multiple dynamic instances of the same load may be issued before
an earlier instance updates the predictor. Such lookups therefore
use the older LAT and SS state. The implementation does not
perform speculative predictor updates to compensate for this
lookup-to-training delay.

\subsection{Value Cache and Load Bypass}
\label{sec:impl-vcache}

V-ASTRA extends ASTRA address prediction with a separate value cache and a guarded load-bypass mechanism. The implementation
has two distinct stages: candidate selection in the Execute stage
and the actual bypass decision in the load-store queue (LSQ).
A load becoming a V-ASTRA candidate does not guarantee that it
will eventually bypass memory.

\subsubsection{Value Cache Organization and Updates}

The value cache is implemented in \texttt{vcache.hh} and
\texttt{vcache.cc}. It is a 4 KiB direct-mapped structure containing
256 entries, each holding 16 bytes. The index and tag are derived
from the virtual address:

\[
    \operatorname{index}(A) =
    (A \gg 4) \mathbin{\&} 255,
\]

\[
    \operatorname{tag}(A) = A \gg 12.
\]

Each entry stores an occupancy bit, a virtual-address tag,
16 data bytes, and a 16-bit validity mask. Each bit in the
mask represents the validity of one byte. A lookup is considered
a hit only when the entry is occupied, the tag matches, and
all bytes requested by the load are valid.

The value cache is updated using completed memory operations through
\texttt{Execute::vcacheRetire()}. Loads allocate an entry if
the indexed entry is unoccupied, replace an entry when its tag
differs, or update an existing entry when the tag matches.
Only the bytes actually read by the load are copied into the
cache and marked valid. The implementation does not fill the
neighboring bytes of a 16-byte entry.

Stores update an existing entry only when its tag matches.
A store to an address not currently represented in the value cache does
not allocate a new entry. Instead, it is recorded as a store miss.

Special operations and uncacheable accesses invalidate the
affected bytes of matching entries. Special operations include
atomics, load-linked/store-conditional operations, prefetches,
micro-operations, swaps, cache-maintenance operations, masked
accesses, and unsupported memory accesses. Uncacheable and
strictly ordered accesses are also handled conservatively.

Loads crossing a 16-byte boundary are neither looked up nor
recorded in the value cache. Stores crossing such a boundary invalidate
the affected bytes in both entries.

The value-cache update path is invoked for completed operations that
have not faulted or been discarded. A load completed through
V-ASTRA bypass does not update the value cache, since it did not obtain
its data through the ordinary memory-access path. The implementation also flushes the value cache after a committed
\texttt{ecall}.

\subsubsection{V-ASTRA Candidate Selection}

Candidate selection is performed by
\texttt{Execute::vastraDecide()} in
\texttt{executeMemRefInst()}. It is invoked after
\texttt{astraCaptureBase()} and before the ordinary memory
access is initiated.

The first step is to validate the predicted base. The value in
\texttt{astraPredBase} is compared with the actual captured base
in \texttt{astraActualBase}. If the base is unavailable or
the comparison fails, the load is rejected as a candidate.

When the base is correct, the effective address is calculated
using the predicted base and the instruction's decoded
displacement \(D\), a signed 12-bit value:

\[
    EA = B + D.
\]

The effective-address helper accepts
uncompressed RISC-V load instructions and rejects unsupported
instruction formats, including compressed loads.

Before accepting the candidate, the implementation checks
for older stores that may overlap the predicted address.
Issued stores, atomics, and store-conditionals are tracked
in a per-thread store list. Their memory-request addresses,
when known, are recorded through the LSQ request path.

The candidate-selection stage scans this list using an
8-byte range, which conservatively covers the supported
load sizes before the actual request size is available.
An older store with an unknown address causes a rejection,
while an overlapping store produces a store-conflict rejection.
A known conflict takes precedence over an unknown address.
Because memory instructions enter the LSQ in program order, an older
store whose address is still unknown is not expected at this point.
The check is retained as a conservative guard.

If the base is correct and no older store conflict or unknown
address is found, the load is marked as a V-ASTRA candidate.
The implementation saves the predicted effective address
and a snapshot of the corresponding value-cache entry, including its
tag-hit status, byte-validity information, and data.

Importantly, a candidate does not require a value-cache hit at this
stage. The candidate decision records the information needed
for the later bypass decision. It does not yet establish that
the value cache contains valid data for the requested bytes.

\subsubsection{LSQ Bypass and Safety Checks}

The actual bypass is attempted later in
\texttt{LSQ::tryToSendToTransfers()}, when the load request
is ready to move from the LSQ request queue to the transfers
queue. This point is reached after address translation, so
the TLB access has already occurred.

The LSQ first enforces its ordinary request-handling conditions.
The request must be ready for transfer, belong to the current
instruction stream, have no translation fault, and have space
in the transfers queue. The request must be a load, with no
older store remaining in the transfers queue. It must also
not require load-linked/store-conditional handling. The store
buffer must not provide full or partial address-range coverage:
full coverage is handled by ordinary store-to-load forwarding,
while partial coverage requires the load to wait.

The LSQ then invokes \texttt{Execute::vastraTryBypass()}.
This function checks the candidate and the actual memory
request before supplying data from the value cache. The checks include:

\begin{enumerate}
    \item V-ASTRA bypass must be enabled, and the request
    must not already have fallen back to ordinary memory access.

    \item The load must have been decided as a V-ASTRA
    candidate.

    \item The request must be a supported single load access:
    a read request with valid packet data and without unsupported
    properties such as atomic, prefetch, micro-operation,
    load-linked/store-conditional, swap, cache-maintenance, masked,
    uncacheable, strictly ordered, or local access types.

    \item The request must fit entirely within a 16-byte
    value-cache entry.

    \item The actual virtual address must match the effective
    address saved when the candidate was selected.

    \item Older stores are checked again using the actual
    request size.

    \item The candidate's value-cache snapshot must have a matching
    tag and all requested bytes must be valid.

    \item A fresh value-cache lookup must still contain the same
    requested bytes as the original candidate snapshot.
\end{enumerate}

The store scan is repeated at the LSQ send point using the
actual request size. The earlier decision used an 8-byte
range because the load size was not yet available. The second
scan provides a check using the precise access size.

The fresh value-cache lookup is also important because the cache may
have changed between candidate selection and the LSQ send
point. The implementation requires the requested bytes to
remain valid and equal to those in the original snapshot.
The LSQ's store-related gating and this fresh lookup together
protect against changes occurring between the two stages.

If any check fails, the load falls back to its ordinary
memory-access path. The implementation records the applicable
fallback reason and marks the request as having fallen back,
preventing repeated attempts from counting the same fallback
more than once.

If all checks succeed, the fresh value-cache bytes are copied into
the load's response packet and \texttt{makeResponse()} is
called. The instruction is marked as bypassed, and the LSQ
treats the request as completed. The ordinary memory request
is suppressed: no packet is sent to memory and the corresponding
memory-access issue counters are not incremented. Address
translation has already taken place, so the bypass removes the
memory access but not the translation.

The bypassed load subsequently follows the normal completion
and in-order commit path. When it reaches the head of the
in-flight queue, \texttt{completeAcc()} processes its response
data and writes it to the architectural destination register. The load therefore still completes
as an ordinary instruction from the perspective of architectural
retirement, even though its data was supplied by a synthetic
response.

A successful V-ASTRA bypass does not itself update the value cache.
The bypassed-load state is recognized by the retirement logic,
which records the bypass outcome and skips the ordinary value-cache
update.

\subsubsection{Strict Verification and Configuration}

The implementation also provides a strict verification mode
through \texttt{enableVAstraVerify}. In this mode, a candidate
that has a matching tag and valid requested bytes is allowed
to execute through the ordinary memory path. When its real
load data becomes available, the saved value-cache snapshot is compared
against the bytes returned by memory.

A matching value increments the verified-hit statistic.
A mismatch produces a panic containing information about the
instruction, predicted base, actual base, effective address,
and predicted and actual data. Verification therefore checks
whether a candidate's saved value-cache data agrees with the result
obtained through the normal memory-access path. It does not
modify the loaded value.

Strict verification cannot be enabled together with actual
V-ASTRA bypass. A bypassed load does not receive an ordinary
memory response, so there is no independent memory value
against which its value-cache data can be compared.

Value-cache use is restricted to syscall-emulation mode, a simulation
with a single CPU object, and a single hardware thread.
These restrictions are enforced during CPU initialization.
The implementation does not detect memory writes performed
by external agents such as DMA.

The relevant configuration parameters are
\texttt{enable\allowbreak VCache},
\texttt{enable\allowbreak VAstra\allowbreak Diag},
\texttt{enable\allowbreak VAstra\allowbreak Verify}, and
\texttt{enable\allowbreak VAstra\allowbreak Bypass}. All are
disabled by default.

Separate counters distinguish candidate selection from actual bypass
completion. The difference is most visible in the V-cache-miss
workload, which produced 800 candidates but no bypasses
(Section~\ref{subsec:standalone-predictor-performance}).

\subsection{Scoreboard Release and Dependent-Load Forwarding}
\label{sec:impl-sb}

MinorCPU uses a scoreboard to track the readiness of source
registers. The scoreboard maintains timing information for
outstanding register writes, including an estimated return cycle
and the instruction sequence number of the latest writer. It
does not store register values. Operands are read from the
architectural register file when an instruction executes through
\texttt{commitInst} or initiates its memory access, except where
the LVP operand substitution or the forwarding override described
below is active.

By default, MinorCPU assumes a memory latency through
\texttt{extraAssumedLat}. With the default value of 2, load
destinations are not marked unpredictable, and dependent
instructions may issue before the load commits. The dependent
operation can nevertheless remain delayed until its producer
commits. To evaluate an explicit stall-on-use policy, the
\texttt{vp\_stall\_on\_use.py} configuration sets
\texttt{extraAssumedLat} to zero. In this configuration, a load's
destination is marked unpredictable, preventing dependent
instructions from issuing until that mark is cleared.

\subsubsection{Design~A: Scoreboard Release}

Design~A allows dependents of a successfully bypassed load to
become eligible for issue before the producer commits. When a
load is issued under stall-on-use timing, its destination
registers are marked unpredictable in the scoreboard. The
implementation records this condition in the dynamic instruction
using \texttt{sbUnpredictable}.

The unpredictable mark is released only after the V-ASTRA bypass
has succeeded and the synthetic load response has been created.
At this point, all bypass checks have passed and the load's
result is known. The implementation calls
\texttt{releaseUnpredictable} for the producer's destination
registers, decrementing their unpredictable-result counts.
It does not modify the ordinary result counts, the recorded
writer sequence numbers, the return-cycle information, or the
architectural register file.

The instruction records that its scoreboard mark has been
released through \texttt{vaScore\allowbreak Released}. During subsequent
commit or discard cleanup, this flag prevents the unpredictable
count from being decremented a second time. If the bypass
attempt fails, no release takes place, and the load follows the
ordinary memory-access path.

This mechanism changes issue readiness but does not forward the
producer's value. A dependent instruction still reads its
operand from the architectural register file when it executes.
Consequently, it can enter its functional unit earlier, but its
execution continues to use the value made available by the
producer's eventual commit. In particular, the release does not
change the scoreboard's recorded writer sequence number.
Therefore, a dependent load still cannot initiate its memory
access early through Design~A alone.

Under the default memory-latency assumption, loads are not marked
unpredictable, so the release has no effect.

\subsubsection{Design~B: Dependent-Load Forwarding}

Design~B extends dependency release by forwarding the validated
producer value to a dependent load. It maintains a per-thread
forwarding-record map indexed by the producer's instruction
sequence number. Each record contains the producer's destination
register, stream identifier, value, and forwarding state.

A record is created when an eligible V-ASTRA load is identified
as a candidate, provided its base is not a provisional LVP-derived
value (Section~\ref{sec:impl-provisional}). The initial state is \texttt{Pending}, since
candidate identification alone does not establish that the
bypass will succeed. The record becomes \texttt{Validated} only
after the final bypass succeeds. If the candidate falls back to
ordinary memory access, is abandoned, or is found unsuitable for
forwarding, its record is marked \texttt{Failed}. Some candidates
that do not reach the final bypass checks, such as those affected
by a translation fault, can remain pending until they are
removed during instruction cleanup.

At publication, the forwarded value is reconstructed from the
freshly checked value-cache bytes. The implementation supports
the integer load operations \texttt{lb}, \texttt{lbu},
\texttt{lh}, \texttt{lhu}, \texttt{lw}, \texttt{lwu}, and
\texttt{ld}. It assembles the requested bytes in little-endian
order and applies the appropriate sign or zero extension.
Floating-point loads and other unsupported load forms are not
eligible producers.

A dependent load can use forwarding only when it is an ordinary
load with exactly one integer source and is at the front of the
memory functional unit. The forwarding record must correspond to the producer
identified by the consumer's recorded
\texttt{instToWaitFor}. The producer and consumer must belong
to the same stream, and the producer's destination register must
match the consumer's source register. The record must also be in
the \texttt{Validated} state. A pending or failed record does not
admit the consumer. The consumer remains subject to the ordinary
dependency check and is reconsidered in subsequent cycles.

For an admitted consumer, the implementation temporarily
substitutes the forwarded value when the instruction reads its
source operand. This substitution is active only around
\texttt{commitInst}, through a scoped operand-override mechanism
implemented in \texttt{ExecContext::\allowbreak getRegOperand}. The
forwarded value is therefore used to calculate the dependent
load's effective address, without writing it into the
architectural register file. The consumer can then proceed to
the load-access path before its producer commits.

The forwarded value is also used when ASTRA captures the
consumer's base register. Consequently, a dependent load that
uses forwarding can itself become an eligible V-ASTRA candidate,
bypass successfully, and publish a value for a subsequent
dependent load. This permits chains of forwarded loads, subject
to the same admission and bypass conditions.

Once an admitted consumer enters the LSQ, it is
marked as accepted. When its memory response is handled, the
implementation checks that the architectural source register
agrees with the value previously forwarded. A mismatch triggers
a panic rather than speculative recovery. This check is possible
because the older producer has committed by the time the
dependent load completes. During instruction cleanup, the
producer's forwarding record is erased, and an accepted
consumer is classified as committed, squashed, or faulted.
Drain checks ensure that no forwarding records or live
forwarded consumers remain.

Forwarding requires V-ASTRA bypass, scoreboard release, and
early memory issue to be enabled. It is disabled by default and
requires a single hardware thread, as do all value-cache-based
mechanisms. The experiments use the \texttt{SingleThreaded} commit
policy. Unlike LVP, forwarding has no dedicated wake-up path for the
other thread policies. Forwarding is also incompatible with LVP
address speculation unless META is enabled.

Neither design provides general register forwarding to arithmetic
instructions, and neither constitutes out-of-order execution.

\subsection{META Integration}
\label{sec:impl-meta}

META is integrated into MinorCPU's early memory-issue path. It
arbitrates between the value supplied by V-ASTRA and an LVP
prediction when a dependent load attempts to issue before its
producer has committed. The selection is performed only for
eligible dependent loads at the front of the memory functional
unit. META re-evaluates the selection on each commit pass until
the consumer is admitted or can proceed through ordinary
dependency handling.

This differs from the paper's META, which selects one prediction
for each load at issue time. Here, the selection concerns the source
of a dependent load's base operand. A selection only permits the
consumer to issue early in that cycle. Its source is recorded only if
its memory request actually enters the LSQ. A consumer
that is evaluated but never admitted early is classified as having
used no prediction source.

\subsubsection{Selection Policy}

The selection logic first identifies the producer associated
with the dependent load. For V-ASTRA, it checks the forwarding
record corresponding to the producer's instruction sequence
number, including the stream and destination-register match.
The record's state determines whether the producer's value is
available for use.

The selection follows this priority:

\begin{enumerate}
    \item \textbf{Validated V-ASTRA value:} If the producer has
    a validated forwarding record and the consumer satisfies
    the forwarding-admission conditions, META selects V-ASTRA.

    \item \textbf{Pending V-ASTRA outcome:} If the producer's
    forwarding record is still pending, the consumer waits.
    This remains the case even if the LVP admission conditions
    are already satisfied. META does not use LVP while the
    V-ASTRA outcome is undecided.

    \item \textbf{LVP prediction:} If the V-ASTRA record is
    failed or absent, META considers LVP. The existing LVP
    admission conditions must hold, including a matching
    speculative producer, an outstanding predicted load, and
    a matching stream.

    \item \textbf{No prediction:} If neither source can admit
    the consumer, it waits for the ordinary dependency
    condition to be satisfied.
\end{enumerate}

The implementation uses the existing V-ASTRA forwarding
admission check and LVP speculative-address admission check.
META does not replace either mechanism's individual eligibility
rules. Instead, it determines which mechanism, if any, is
allowed to serve the consumer during that evaluation.

The implementation permits both admission checks to succeed in the
same evaluation, although this case was not observed in the retained
directed tests. Only one source is then selected. V-ASTRA takes
priority over LVP, and the implementation checks that the two acceptance paths
cannot both be active for the same instruction.

\subsubsection{Source Attribution and Recovery}

The selected source is recorded when the consumer is accepted
into the LSQ. A consumer admitted through LVP is
marked as LVP-served, and its producer's live-consumer count
is incremented. A consumer admitted through V-ASTRA is marked
as V-ASTRA-served and does not contribute to the LVP
producer's live-consumer count.

This distinction is important because the two sources have
different recovery requirements. An LVP-served consumer uses
a speculative value, so an incorrect producer prediction
requires the usual LVP recovery. Recovery itself is triggered by
the LVP rule: an incorrect prediction that an accepted consumer
used. META additionally counts the live LVP consumers of each
producer. When the producer's memory response is processed, it stops
the simulation if such consumers exist but the required recovery
has not been raised. This count is therefore a consistency check
on the attribution, not the recovery trigger.

A V-ASTRA-served consumer does not mark the producer as having
a live LVP user. Therefore, even if the producer's independent
LVP prediction is later found to be incorrect, that prediction
does not cause recovery on behalf of the V-ASTRA-served
consumer. The LVP prediction is still validated and updated
through the normal LVP path.

At instruction removal, META classifies the consumer according
to its recorded source and terminal outcome: committed,
squashed, or faulted. Consistency checks ensure that the
recorded source agrees with the corresponding LVP or V-ASTRA
consumer state, and that an instruction has not been accepted
through both mechanisms.

\subsubsection{Configuration}

META is disabled by default. Enabling it requires LVP,
LVP address speculation, and V-ASTRA forwarding. Through the
existing configuration checks, this also requires the ASTRA
predictor, value cache, diagnostic support, V-ASTRA bypass,
scoreboard release, and early memory issue. The test-only
options that skip syscall flushing or store-conflict scanning
are incompatible with META. META permits LVP address speculation
and V-ASTRA forwarding to be enabled together, lifting the earlier
incompatibility between these mechanisms.

\subsection{Experimental ASTRA Interaction for LVP-Served Loads}
\label{sec:impl-provisional}

The default META configuration does not use an LVP-supplied operand
as ASTRA base information. Consequently, a load whose base address
is obtained through LVP cannot contribute to subsequent ASTRA
training or become a V-ASTRA candidate. In this default \texttt{off}
mode, \texttt{astraCaptureBase()} does not capture an LVP-supplied
base, so the load has no valid ASTRA base metadata and does not train
ASTRA. This limits the interaction between the two prediction
mechanisms, particularly for workloads in which LVP serves many
dependent loads.

To investigate this limitation, the implementation introduces the
optional \texttt{metaAstra\allowbreak LvpLoads} parameter. It controls how a
load using an LVP-predicted base contributes information to ASTRA.
Three modes are supported: \texttt{off}, \texttt{train}, and
\texttt{candidate}. The experiment distinguishes between recording
a speculative base for future predictor training and allowing that
base to participate immediately in V-ASTRA candidate selection.

The three modes modify the treatment of LVP-derived bases without
changing META's source-selection policy. They are implemented
through the base-capture path and the associated provisional-state
handling in the Execute stage.

\subsubsection{Provisional Base Capture and Validation}

ASTRA normally captures a load's actual base-register value when
the memory instruction is processed by
\texttt{executeMemRefInst()}, immediately before its memory access
is initiated. When the base operand is obtained through ordinary
execution, the implementation records the architectural register
value. A forwarded V-ASTRA operand is used instead when the
corresponding forwarding scope is active.

An LVP-served load requires different treatment. Its base operand
is supplied by a speculative prediction, which has not yet been
validated against the producer's actual result. The implementation
therefore distinguishes between two properties of the captured
base:

\begin{itemize}
    \item \texttt{astraBaseProvisional} indicates that the base was
    obtained from an unvalidated LVP operand.

    \item \texttt{astraBaseValid} determines whether the captured
    base may be used by the ASTRA/V-ASTRA decision and the
    subsequent training path.
\end{itemize}

Training also requires the load to commit without a fault and
without an unresolved provisional base. A load that commits while
its base is still provisional is treated as an internal error.

These flags are independent. In particular, candidate mode allows
a provisional base to be used for candidate selection by setting
\texttt{astraBase\allowbreak Valid}, while retaining
\texttt{astraBase\allowbreak Provisional} to indicate that the underlying LVP
operand has not yet been validated. Thus, being valid for the
candidate decision does not mean that the base is already known
to be architecturally correct.

For the \texttt{train} and \texttt{candidate} modes, the captured
provisional base is the value supplied through the active LVP
operand substitution. The implementation records this value and
increments the provisional-base creation statistic. The mode
determines whether it is immediately available to the
V-ASTRA candidate-selection stage.

A provisional base becomes validated only when the load's own
memory response is handled. Before resolving the provisional
state, the implementation checks that the architectural source
register agrees with the LVP value used by the load. If the check
succeeds, the base is replaced with the architectural register
value, \texttt{astraBaseProvisional} is cleared, and
\texttt{astraBaseValid} is set. This is the point at which the
provisional base becomes suitable for normal ASTRA training.
This resolution takes place at the start of response handling,
before the response's own fault, if any, is processed. A load whose
LVP operand was correct but whose response carries a translation
fault is therefore counted as validated, and then retires as faulted
without training ASTRA.

The actual ASTRA update still occurs through the ordinary
retirement path. A correctly resolved base is therefore not
immediately written into the predictor tables when the memory
response arrives. Instead, the load trains ASTRA when it
subsequently retires successfully.

If the load is discarded, or is re-executed after an early fault,
before successful resolution, its provisional state is removed
rather than being used for training. The implementation also enforces an invariant at
instruction removal: a load must not commit while its base remains
provisional. A committed instruction with unresolved provisional
state is treated as an internal error.

\subsubsection{Mode \texttt{train}: Deferred ASTRA Training}

The \texttt{train} mode allows an LVP-served load to contribute
base information to ASTRA, but only after the LVP operand has
been validated.

When the load is processed with an active LVP operand, the
implementation captures the predicted base and marks it as
provisional. However, it leaves \texttt{astraBaseValid} false.
The subsequent V-ASTRA decision therefore rejects the load due
to missing valid base metadata. The provisional value does not
participate in candidate selection, and the load cannot use this
path to obtain a value-cache bypass.

When the load's own response is handled, the implementation
checks the architectural source register against the value
previously supplied by LVP. If the values agree, the provisional
base is resolved and replaced with the architectural value.
The load can then participate in the normal ASTRA retirement
update.

Training is consequently deferred until successful retirement.
The experimental mode does not train ASTRA directly from an
unvalidated prediction. Instead, it preserves the opportunity to
train using the corresponding architectural base once the LVP
operand has been confirmed.

If the LVP prediction is incorrect and has been used by the load,
the normal LVP recovery mechanism redirects execution. The
younger load is discarded, its provisional base is dropped, and
ASTRA does not train using that speculative observation. This
ensures that a wrong LVP operand does not become a permanent
ASTRA training value.

The \texttt{train} mode therefore increases the information
available to ASTRA without allowing an unvalidated base to
affect the current V-ASTRA candidate decision.

\subsubsection{Mode \texttt{candidate}: Provisional V-ASTRA Candidacy}

The \texttt{candidate} mode is more aggressive. It permits the
captured LVP operand to participate immediately in V-ASTRA
candidate selection, even though the operand has not yet been
validated.

For an LVP-served load, the implementation captures the
predicted base and sets both
\texttt{astraBase\allowbreak Provisional} and
\texttt{astraBase\allowbreak Valid}. The V-ASTRA decision can therefore
compare ASTRA's predicted base with the LVP-supplied base.

If the two bases match, the effective address is formed using
the decoded load displacement:

\[
    EA = B + D,
\]

where \(B\) is the base predicted by ASTRA and \(D\) is the
signed 12-bit displacement encoded in the load instruction. Since the provisional base is the LVP operand,
this comparison establishes agreement between the ASTRA
prediction and the LVP-supplied value. It does not independently
establish that the LVP value is architecturally correct.

The candidate must still pass the normal V-ASTRA safety checks.
These include the initial older-store scan, the value-cache
snapshot, and the subsequent LSQ-stage checks. At the LSQ send
point, the implementation verifies the actual request address,
repeats the store-conflict scan using the actual access size,
checks the requested bytes in the saved value-cache snapshot,
and performs a fresh value-cache lookup to ensure that the
requested bytes still match.

The verified request address is the address that
\texttt{initiateAcc()} produced from the LVP operand. Its agreement
with the saved candidate address therefore confirms consistency with
ASTRA's prediction, not architectural correctness.

Thus, candidate mode does not bypass the ordinary V-ASTRA
store-safety or value-cache consistency requirements. Its
additional speculation concerns the source of the base address:
the LVP operand may be used before its correctness is established.
This is acceptable because the load has already been accepted as an
LVP consumer: its result exists only as a pending response in the
LSQ, and if the LVP prediction proves incorrect, the
normal LVP recovery discards the load before it can commit.

A successful bypass can therefore occur before the LVP
prediction has been validated. The load receives a synthetic
response containing the value-cache data and subsequently
follows the normal in-order completion path. If the LVP operand
is later confirmed, the provisional base is resolved and the
load can train ASTRA at retirement. As with other bypassed
loads, the successful bypass does not update the value cache
through the ordinary memory-response update path.

A provisional load is not allowed to create a dependent-load
forwarding record. A successful candidate decision only shows that
ASTRA's predicted base agrees with the LVP operand, not that the
operand is correct. Publishing the value would therefore let a younger
dependent load use data derived from an unvalidated base. Such
dependents wait for the producer through ordinary dependency handling.

Correctness therefore relies on the existing LVP recovery. If the
original producer's actual value differs from the prediction, the
recovery redirect discards the younger candidate-mode load before it
commits. The forced-misprediction tests verify the following order
from the simulation trace for every such load, where each event occurs
in the same cycle as the previous one or later:
\[
    \text{bypass}
    \;\rightarrow\;
    \text{validation}
    \;\rightarrow\;
    \text{redirect}
    \;\rightarrow\;
    \text{discard}.
\]

\subsubsection{Fault Handling and Provisional-State Retry}

The provisional-base mechanism also requires careful handling
when an early-issued memory instruction faults before entering
the LSQ. This case exposed a lifecycle defect in the initial
implementation.

An early-issued load can use an LVP operand to form its effective
address and then encounter an alignment fault during
\texttt{initiateAcc()}. On RISC-V, \texttt{initiateAcc()} raises such
a fault only when alignment checking is enabled. In this situation,
the request has not entered the LSQ. MinorCPU prevents the instruction
from continuing through the early-issue path and allows it to be
executed again when it reaches the head of the in-flight queue.

This case differs from a translation fault. A translation fault is
recorded with a request that has entered the LSQ and is processed
when the load's response is handled. As described in
Section~\ref{sec:impl-provisional}, the provisional base is resolved
at the start of that response handling, before the fault itself is
processed.

On re-execution, the temporary LVP operand-substitution scope
from the earlier attempt is no longer active. The instruction
therefore captures its base from the ordinary architectural
register state. However, before the fix, the provisional flag
created during the earlier attempt could remain set. The
re-executed instruction could consequently reach retirement
with a stale provisional state, triggering the invariant that
a committed load must not retain an unresolved provisional
base.

The fix handles this at the beginning of every new ASTRA
base-capture attempt. If a provisional base remains from an
earlier attempt, it is explicitly dropped before the new base
is captured. The base-valid flag is cleared, and the provisional
drop statistic is updated. The current execution attempt can
then establish its own base state without inheriting the
provisional flag from the failed attempt.

The candidate-mode retry path has an additional subtlety:
the first V-ASTRA decision is retained across the early fault.
On re-execution, the ordinary architectural base may differ
from the address used by the original provisional decision.
The later bypass check rejects the request when its actual
address does not match the saved candidate address. If the LVP
operand had been correct, the re-executed request would use the
same misaligned address and fault again, so no request would reach
the bypass point.

\subsubsection{Configuration}

The \texttt{metaAstraLvpLoads} parameter accepts the values
\texttt{off}, \texttt{train}, and \texttt{candidate}, with
\texttt{off} as the default. Selecting either experimental mode
requires META to be enabled. Statistics under
\texttt{system.cpu.meta} record how many provisional bases are
created, validated, and dropped. Validation and performance results
for the three modes are given in
Sections~\ref{subsec:meta-integration-recovery-validation}
and~\ref{subsec:provisional-base-performance}.
